% notes/v3/h1_constant.tex  --  the H^1 leak constant of the printed method (item 5)
% Fragment in the macros of paper/main.tex; labels hc:.  Scripts: code/h1_limit.py; logs: logs/h1_limit_*.log

\section{The $H^1$ leak constant}\label{hc:sec}

Corollary~\ref{cor:Bp} identifies the slip and energy constants of $\GS$ (both $\to1$), but not the $H^1$ constant
$\norm{\nabla(\ut-u_h)}\,\mu/(hG)$, because the remainder $r_h=\delta_R-(h/\mu)S_h$ of Theorem~\ref{thm:C} has the same
$H^1$ order as $(h/\mu)S_h$. We identify the limit of $(\mu/h)\delta_R$: it is the Stokes flow in $\Omega$ driven by the
leak as \emph{Dirichlet} data. The $H^1$ constant is its Dirichlet energy, a global quantity that depends on the outer
box. For test P it equals $2.7565157\ldots$

\subsection{The limit problem}

\begin{definition}[Leak flow]\label{hc:def}
Let $(U,P)$ solve
\[
-\Delta U+\nabla P=0,\quad \div U=0\ \text{in }\Omega,\qquad U=(p-\pbar)\,e_r=-(p-\pbar)\,n\ \text{on }\Gamma,\qquad U=0\ \text{on }\partial R ,
\]
and put $K:=\norm{\nabla U}_{L^2(\Omega)}/G$.
\end{definition}

The data are compatible, since $\int_\Gamma U\cdot n=-\int_\Gamma(p-\pbar)=0$. Hence $(U,P)$ exists and is unique ($P$ up to
a constant), $U=\curl\Psi$ with a single-valued stream function, and $U|_\Gamma=w|_\Gamma$ with $w=\curl\varphi$ the test field of
Section~\ref{sec:printed}. $U$ is smooth up to $\Gamma$ (smooth boundary and data) and $U\in H^2(\Omega)$ (the corners of $R$
are convex). Since $-\Delta U=-\nabla P$, for every divergence-free $V$ with the same boundary values
$\int\nabla U:\nabla(V-U)=-\int\nabla P\cdot(V-U)=0$, so
\begin{equation}\label{hc:min}
\norm{\nabla U}^2=\min\Bigl\{\norm{\nabla V}^2_{L^2(\Omega)}:\ V\in H^1(\Omega)^2,\ \div V=0,\ V|_\Gamma=(p-\pbar)e_r,\ V|_{\partial R}=0\Bigr\}.
\end{equation}
In particular $K$ depends on $p|_\Gamma$ \emph{and} on the outer boundary: it is not a local quantity.

\subsection{Convergence}

Extend $\Psi$ smoothly across $\Gamma$ to a fixed neighbourhood and put $\tilde U:=\curl\tilde\Psi$ on $\Oh$, so $\div\tilde U=0$;
extend $P$ smoothly to $\tilde P$. Then $F_U:=-\Delta\tilde U+\nabla\tilde P$ vanishes on $\Omega$, is bounded, and is supported in
the strip $S_h$.

\begin{lemma}\label{hc:approx}
There is $z_U\in\Zh$ with $\tnorm{\tilde U-z_U}\le E_U:=C\bigl(h+(1+(\gamma\hG)^{1/2})h^k\bigr)$.
\end{lemma}

\begin{proof}
Repeat the proof of Lemma~\ref{lem:approx} with $\ut$ replaced by $\tilde U$ and $g=0$. Away from $\Gamma$ only $\tilde U\in
H^2$ is available, so the interpolation step gives $\norm{\nabla(\tilde U-I_h\tilde U)}\le Ch\norm{U}_{H^2}$. Near $\Gamma$, $\tilde U$
is smooth and the boundary terms are as before. The net flux $m=\int_{\Gh}(I_h\tilde U-\tilde U)\cdot n$ is $O(\hG^{\sigma_k})$
because $\int_{\Gh}\tilde U\cdot n=\int_{\Oh}\div\tilde U=0$; the bubble and the Bogovskii correction (Lemma~\ref{lem:infsup})
are unchanged.
\end{proof}

\begin{theorem}[The leak flow is the $H^1$ limit]\label{hc:thm}
Let $\gamma\ge\gamma_0$ and let $\delta_R\in\Zh$ solve $N_h(\delta_R,v)=\Rc(v)$ for all $v\in\Zh$ (Theorem~\ref{thm:C}, step 1).
Then
\[
\Bigl\|\nabla\Bigl(\frac\mu h\delta_R-\tilde U\Bigr)\Bigr\|_{\Oh}
\le\tnorm{\frac\mu h\delta_R-z_U}+E_U
\le C_U\Bigl[\min\bigl(\gamma\hG^{1/2},(\gamma\hG)^{1/2}\bigr)+\hG^{1/2}+(h/\mu)^{1/2}+E_U\Bigr],
\]
where $C_U$ depends on (M0)--(M2), $k$, $L$ and on norms of $p$ and $U$ near $\Gamma$.
\end{theorem}

\begin{proof}
Put $\epsilon:=\delta_R-(h/\mu)z_U\in\Zh$. For $v\in\Zh$,
\begin{equation}\label{hc:erreq}
N_h(\epsilon,v)=\Rc(v)-\frac h\mu N_h(\tilde U,v)+\frac h\mu N_h(\tilde U-z_U,v).
\end{equation}
\begin{enumerate}
\item \emph{Green's formula for $\tilde U$.} As in Lemma~\ref{lem:erroreq}, $\div D(\tilde U)=\Delta\tilde U$, $\div v=0$, $v|_{\partial R}=0$ and
$\int_{\Gh}v\cdot n=0$ give, for any constant $c$,
\[
N_h(\tilde U,v)=(F_U,v)+\ip{\sigma_U}{v}-\ip{\tilde U}{\dn v}+\frac\mu h\ip{\tilde U}{v},\qquad
\sigma_U:=(\nabla\tilde U)^{\mathsf T}n-(\tilde P-c)\,n .
\]
\item \emph{The penalty terms cancel to leading order.} With $\Rc(v)=-\ip{(\pt-\pbar)n}{v}$, \eqref{hc:erreq} becomes
\[
N_h(\epsilon,v)=-\ip{\zeta}{v}-\frac h\mu\Bigl[(F_U,v)+\ip{\sigma_U}{v}-\ip{\tilde U}{\dn v}\Bigr]+\frac h\mu N_h(\tilde U-z_U,v),
\qquad \zeta:=(\pt-\pbar)n+\tilde U .
\]
\item \emph{The mismatch $\zeta$.} On $e\in\EG$, by Lemma~\ref{lem:chords} and $\tilde U(x^\ast)=-(p-\pbar)(x^\ast)n(x^\ast)$,
\[
\zeta=(p-\pbar)(x^\ast)\bigl(n_e-n(x^\ast)\bigr)+O(\hG^2)=(p-\pbar)(m_e^\ast)\,s\,\tau_e+O(\hG^2),
\]
exactly the functional $m$ of Theorem~\ref{thm:C}, step~4 (without the $h^k$ term). Hence
$|\ip{\zeta}{v}|\le C\hG^{3/2}\norm{\nabla v}$. Alternatively $\norm{\zeta}_{\Gh}\le C\hG$ gives
$|\ip{\zeta}{v}|\le C\hG(h/\mu)^{1/2}\tnorm v$.
\item \emph{The $O(h/\mu)$ terms.} $|(F_U,v)|\le C\hG^2\norm{\nabla v}$ (Lemma~\ref{lem:ext}); $|\ip{\sigma_U}{v}|\le
C\norm{v}_{\Gh}\le C(h/\mu)^{1/2}\tnorm v$. On $\Gh$, $\tilde U$ is edgewise smooth, continuous at the vertices, and
$\tilde U\cdot\tau_e=-(p-\pbar)(x^\ast)\,n(x^\ast)\cdot\tau_e+O(\hG^2)=O(\hG)$. So Lemma~\ref{lem:cancel} applies with
$S=\tilde U$: $|\ip{\tilde U}{\dn v}|\le C(\norm v_{\Gh}+\hG^{1/2}\norm{\nabla v})\le C((h/\mu)^{1/2}+\hG^{1/2})\tnorm v$.
Finally $|N_h(\tilde U-z_U,v)|\le ME_U\tnorm v$ (Lemmas~\ref{lem:coercive} and~\ref{hc:approx}).
\item \emph{Close.} Test with $v=\epsilon$ and use Lemma~\ref{lem:coercive}:
\[
c_1\tnorm{\epsilon}\le C\Bigl[\min\bigl(\hG^{3/2},\hG(h/\mu)^{1/2}\bigr)+\frac h\mu\bigl(\hG^2+(h/\mu)^{1/2}+\hG^{1/2}+E_U\bigr)\Bigr].
\]
Divide by $h/\mu$, and use $\hG^{3/2}\mu/h=\gamma\hG^{1/2}$ and $\hG(\mu/h)^{1/2}=(\gamma\hG)^{1/2}$. The first inequality of the
theorem is the triangle inequality with Lemma~\ref{hc:approx}.\qedhere
\end{enumerate}
\end{proof}

\begin{corollary}[The $H^1$ constant]\label{hc:cor}
\begin{enumerate}[label=(\roman*)]
\item For $\GS$ with general data, in the notation of Theorem~\ref{thm:C},
\[
\Bigl\|\nabla\Bigl(\frac\mu h(\ut-u_h)-\tilde U\Bigr)\Bigr\|_{\Oh}\le\text{(bound of Theorem~\ref{hc:thm})}
+C\frac\mu h\Bigl[(1+\gamma^{1/2})\hG^{3/2}+\eps\Bigr].
\]
\item Consequently, if $\hG\to0$ with $h/\mu\to0$, $\gamma\hG\to0$, $\gamma(1+\gamma^{1/2})\hG^{1/2}\to0$ and
$(\mu/h)\eps\to0$ (all of which hold for fixed $\gamma$ and bounded $\rho$), then
\[
\frac{\norm{\nabla(\ut-u_h)}_{\Oh}}{(h/\mu)\,G}\longrightarrow K=\frac{\norm{\nabla U}_{L^2(\Omega)}}{G}.
\]
\item On test P ($u=0$, $f=\nabla p$) one has $\ut=0$, $z=0$ and $\Gc\equiv0$, so $\ut-u_h=\delta_R$ exactly and
Theorem~\ref{hc:thm} applies with no further condition.
\item The remainder of Theorem~\ref{thm:C} converges: $(\mu/h)\,r_h\to U-w$ in $H^1$. The limit is the Stokes correction
of the cut-off field $w$: it vanishes on $\partial\Omega$ and solves $-\Delta(U-w)+\nabla P=\Delta w$. It depends on the
cut-off $\chi_1$, whereas $U$ does not.
\end{enumerate}
\end{corollary}

\begin{proof}
(i) is $\ut-u_h=(\ut-z)+\delta_R+\delta_G$ with steps 2 of Theorem~\ref{thm:C} and Lemma~\ref{lem:approx}. (ii) Multiply out;
$\norm{\nabla\tilde U}^2_{\Oh}-\norm{\nabla U}^2_{\Omega}=\norm{\nabla\tilde U}^2_{S_h}=O(\hG^2)$. (iii) $F=\ft-\nabla\pt=0$ and all terms of
$\Gc$ contain $\ut$. (iv) $r_h=\delta_R-(h/\mu)v_\varphi$ and $\norm{\nabla(v_\varphi-w)}\le Ch^k$.
\end{proof}

\begin{remark}[The Dirichlet-dominated regime]\label{hc:regime}
The hypothesis $\gamma\hG=\mu\hG^2/h\to0$ is the one of Corollary~\ref{cor:Bp}. It fails for the large-$\mu$ runs of
Table~\ref{tab:P} ($\gamma\hG\approx240$ at $N=96$, $\mu=10^4$). There the discrete problem is close to its $\mu\to\infty$ limit,
a discrete Dirichlet problem whose data $-(p-\pbar)n_e$ jumps in direction at every vertex of $\Gh$. The computations below
show that $(\mu/h)\delta_R\to U$ in that regime too, with relative $H^1$ distance decaying roughly like $\hG^{0.8}$, but
Theorem~\ref{hc:thm} does not cover it. Proving convergence uniformly in $\mu\ge\mu_0$ is open.
\end{remark}

\subsection{Rigorous bracket and exact asymptotics}

For $R>1$ let $\mathcal A_R=\{1<r<R\}$ and let $\mathcal E(R)$ be the minimum \eqref{hc:min} over $\mathcal A_R$ (zero data on
$r=R$). Extending by zero shows that the minimum decreases as the domain grows.

\begin{proposition}\label{hc:bracket}
If $\mathcal A_{R_1}\subset\Omega\subset\mathcal A_{R_2}$, then $\mathcal E(R_2)\le\norm{\nabla U}^2\le\mathcal E(R_1)$. In the annulus the
Fourier modes decouple: if $\Psi|_\Gamma=\sum_n(a_n\cos n\theta+b_n\sin n\theta)$, then $\mathcal E(R)=\sum_n(a_n^2+b_n^2)E_n(R)$, where
$E_n(R)=\pi\int_1^R\bigl[f''^2+2n^2((f/r)')^2+(f'/r-n^2f/r^2)^2\bigr]r\,dr$, and $f$ is the combination of
$r^n,r^{-n},r^{n+2},r^{2-n}$ ($r,r^{-1},r^3,r\log r$ for $n=1$) with $f(1)=1$, $f'(1)=0$, $f(R)=f'(R)=0$.
\end{proposition}

For test P, $p-\pbar=-\frac34\sin\theta+\frac14\sin3\theta+\frac12\cos\theta$ on $\Gamma$, so
$\Psi|_\Gamma=\frac34\cos\theta+\frac12\sin\theta-\frac1{12}\cos3\theta$ and $G^2=7\pi/8$. With $R_1=L=2.5$ and $R_2=L\sqrt2$
(40-digit quadrature):
\[
\boxed{\,2.0746004\le K\le3.0525769\,}
\]
The same formulas give the dependence on the box. $E_3(R)\to45\pi$ (the exterior clamped solution is
$f=\frac32r^{-1}-\frac12r^{-3}$). $E_1(R)\to\pi$: the $n=1$ data split into a potential dipole (exterior energy exactly $\pi$) and
a rigid translation, whose energy vanishes like $1/\log R$ (the Stokes paradox). Hence
\[
K(L)\longrightarrow\sqrt{\tfrac{(\frac9{16}+\frac14)\pi+\frac{45\pi}{144}}{7\pi/8}}=\frac3{\sqrt7}\approx1.134\qquad(L\to\infty),
\]
but only logarithmically: the annulus values are $K=1.420$, $1.242$, $1.183$ at $R=10,10^2,10^4$. So the constant
$2.7$--$2.8$ of Table~\ref{tab:P} belongs to the box $L=2.5$, not to the method.

\subsection{The value for test P}

\emph{Series solution.} Write $\Psi=\sum_{n\text{ odd}\le M}f_n(r)\cos n\theta$ for data $\cos\theta$ and $\cos3\theta$ separately (the class of
$\cos n\theta$, $n$ odd, is invariant under both reflections of the square; the $\sin\theta$ datum is the rotation of the
$\cos\theta$ one by $\pi/2$, and cross terms between the two reflection classes vanish). The conditions on $r=1$ are imposed
exactly, mode by mode, and $\Psi=\partial_\nu\Psi=0$ on $\partial R$ by least squares at $16\,000$ Chebyshev-clustered points.
The energy is integrated by Gauss quadrature on eight polar sectors. The value is stable to 9--10 digits in both $M$ and the
quadrature:

\begin{center}
\begin{tabular}{rccc}
\toprule
$M$ & max boundary residual & $\norm{\nabla U}$ & $K$\\
\midrule
15 & $4.4\cdot10^{-3}$ & 4.57012143 & 2.75644148\\
31 & $6.1\cdot10^{-5}$ & 4.57024239 & 2.75651444\\
47 & $5.1\cdot10^{-6}$ & 4.57024449 & 2.75651571\\
63 & $1.1\cdot10^{-6}$ & 4.57024450 & 2.75651571\\
79 & $1.8\cdot10^{-7}$ & 4.57024450 & 2.75651571\\
\bottomrule
\end{tabular}
\end{center}

\[
\boxed{\,K_{\mathrm P}=2.7565157\ (\pm10^{-9})\,}
\]
This lies inside the rigorous bracket.
\cert{code/h1\_limit.py annulus; code/h1\_limit.py mps}{logs/h1\_limit\_series.log}

\emph{Finite element check.} The paper's solver (SuperLU with refinement, as in Section~\ref{sec:numerics}) was run on test P, and
the discrete field was compared \emph{pointwise} with the series $U$ at the quadrature points. Both the normalised $H^1$ error
$K_h$ and the relative distance $d_h:=\norm{\nabla((\mu/h)(\ut-u_h)-\tilde U)}/\norm{\nabla U}$ are shown below. The $\mu=100$ and
$N=96$ values of $K_h$ reproduce Table~\ref{tab:P} to all printed digits.

\begin{center}\small
\begin{tabular}{r|cc|cc|cc|cc}
\toprule
 & \multicolumn{2}{c|}{$\mu=10^2$} & \multicolumn{2}{c|}{$\mu=10^3$} & \multicolumn{2}{c|}{$\mu=10^4$} & \multicolumn{2}{c}{$\mu=10^8$}\\
$N$ & $K_h$ & $d_h$ & $K_h$ & $d_h$ & $K_h$ & $d_h$ & $K_h$ & $d_h$\\
\midrule
16 & 2.4642 & 0.193 & 2.7844 & 0.318 & 3.0562 & 0.539 & 3.3254 & 0.717\\
24 & 2.5629 & 0.147 & 2.7913 & 0.259 & 2.9309 & 0.398 & 3.0452 & 0.497\\
32 & 2.6117 & 0.124 & 2.7880 & 0.222 & 2.8751 & 0.323 & 2.9339 & 0.385\\
48 & 2.6601 & 0.098 & 2.7808 & 0.179 & 2.8283 & 0.246 & 2.8514 & 0.277\\
64 & 2.6843 & 0.084 & 2.7759 & 0.154 & 2.8082 & 0.206 & 2.8208 & 0.226\\
96 & 2.7084 & 0.067 & 2.7701 & 0.124 & 2.7899 & 0.163 & 2.7959 & 0.175\\
\bottomrule
\end{tabular}
\end{center}
\cert{code/h1\_limit.py fem}{logs/h1\_limit\_fem.log}

The measured values $2.71$--$2.79$ bracket the limit $2.7565$ from both sides. The approach runs in two directions.
\begin{itemize}
\item At $\mu=100$, $K_h$ increases towards $K$ and $d_h$ decays like $h^{0.6}$ ($0.193\to0.067$ as $h$ drops by a factor $5.3$).
This is consistent with Theorem~\ref{hc:thm}: for fixed $\mu$, the terms $(h/\mu)^{1/2}$ and $\hG^{1/2}$ are both
$\propto h^{1/2}$, and $(\gamma\hG)^{1/2}\propto h^{1/2}$ too.
\item At large $\mu$ (Remark~\ref{hc:regime}), $K_h$ decreases towards $K$ from above. The excess is the energy of an
element-scale oscillation that is nearly orthogonal to $U$: $K_h^2\approx K^2(1+d_h^2)$, for example $2.7565\sqrt{1+0.175^2}=2.798$
against $2.796$ at $N=96$.
\end{itemize}
At fixed $N=96$ the $\mu$-dependence of Table~\ref{tab:P} ($2.708\to2.770\to2.790$) is the crossover between the two regimes,
not a sign of a different limit.
